\documentclass[11pt]{article}

% ── Packages ─────────────────────────────────────────────────────────
\usepackage[utf8]{inputenc}
\usepackage[T1]{fontenc}
\usepackage{lmodern}
\usepackage[margin=1in]{geometry}
\usepackage{listings}
\usepackage{xcolor}
\usepackage{hyperref}
\usepackage{booktabs}
\usepackage{amsmath,amssymb}
\usepackage{enumitem}
\usepackage{caption}
\usepackage{subcaption}
\usepackage{tabularx}

% ── Listing styles ────────────────────────────────────────────────────
\definecolor{leanKeyword}{RGB}{0,0,180}
\definecolor{leanComment}{RGB}{100,100,100}
\definecolor{leanString}{RGB}{0,120,0}
\definecolor{leanOutput}{RGB}{60,60,60}
\definecolor{codebg}{RGB}{248,248,248}

\lstdefinelanguage{Lean}{
  morekeywords={import,open,namespace,end,def,theorem,example,syntax,
    let,match,return,try,catch,do,if,then,else,fun,where,with,
    partial,private,structure,deriving,in,for,mut,section,
    attribute,initialize,class,instance},
  morekeywords=[2]{use,resolve},
  keywordstyle=\color{leanKeyword}\bfseries,
  keywordstyle=[2]\color{leanKeyword}\bfseries,
  sensitive=true,
  morecomment=[l]{--},
  morecomment=[n]{/-}{-/},
  commentstyle=\color{leanComment}\itshape,
  morestring=[b]",
  stringstyle=\color{leanString},
  basicstyle=\ttfamily\small,
  breaklines=true,
  showstringspaces=false,
}

\lstset{
  language=Lean,
  backgroundcolor=\color{codebg},
  frame=single,
  rulecolor=\color{black!20},
  xleftmargin=2em,
  framexleftmargin=1.5em,
  numbers=left,
  numberstyle=\tiny\color{black!40},
  numbersep=8pt,
  aboveskip=1em,
  belowskip=1em,
}

\lstdefinestyle{output}{
  language={},
  basicstyle=\ttfamily\small\color{leanOutput},
  frame=single,
  rulecolor=\color{black!15},
  backgroundcolor=\color{white},
  numbers=none,
  xleftmargin=2em,
  framexleftmargin=0em,
}

\lstdefinestyle{pseudo}{
  language={},
  basicstyle=\ttfamily\small,
  frame=single,
  rulecolor=\color{black!20},
  backgroundcolor=\color{codebg},
  numbers=none,
  xleftmargin=2em,
  framexleftmargin=0em,
}

\lstdefinestyle{json}{
  language={},
  basicstyle=\ttfamily\small,
  frame=single,
  rulecolor=\color{black!20},
  backgroundcolor=\color{codebg},
  numbers=none,
  xleftmargin=2em,
  framexleftmargin=0em,
}

% ── Review annotation colors ─────────────────────────────────────────
\newcommand{\leo}[1]{\textcolor{purple}{[\textbf{Leo}: #1]}}
\newcommand{\franz}[1]{\textcolor{red}{[\textbf{Franz}: #1]}}
\newcommand{\sam}[1]{\textcolor{orange}{[\textbf{Sam}: #1]}}
\newcommand{\simone}[1]{\textcolor{blue}{[\textbf{Simone}: #1]}}
\newcommand{\moqian}[1]{\textcolor{teal}{[\textbf{Moqian}: #1]}}


% ── Metadata ──────────────────────────────────────────────────────────
\title{\textbf{Content-Addressing Formal Mathematics}}
\author{Xinze Li\textsuperscript{1} \and Marcello Paris \and Simone Severini\textsuperscript{2} \\[4pt]
\textsuperscript{1}University of Toronto \quad \textsuperscript{2}Google}
\date{March 2026}

\begin{document}

\maketitle

\begin{abstract}
We explore content addressing for formal mathematics, introducing a stable identifier determined by mathematical content rather than human chosen names. By computing a cryptographic hash of a canonicalized type expression, we can globally address mathematical facts. We hope this approach helps enable a new paradigm for decentralized knowledge sharing, contribution, and discovery. Building on these addresses, we constructed a decentralized registry where the Lean verifier itself serves as the consensus mechanism. We share a working prototype for Lean 4, evaluate it on standard Mathlib modules, and openly discuss our design space, current limitations, and connections to both content addressing in software engineering and prior work on cross prover interoperability.
\end{abstract}

%\tableofcontents
%\bigskip

% =====================================================================
\section{Introduction}
\label{sec:intro}
% =====================================================================

Software engineers are familiar with content-addressing.  Git identifies
every commit, tree, and blob by the SHA-1 hash of its
content~\cite{git}.  The Nix package manager derives store paths from
the hash of all build inputs~\cite{dolstra2004nix}.  IPFS addresses
files by their cryptographic hash, making retrieval
location-independent~\cite{benet2014ipfs}.  The Unison programming
language names functions by the hash of their syntax
tree~\cite{unison}.

In each case, the core idea is the same: derive the identifier from the
\emph{content itself}, not from a human-chosen name or a filesystem
location.  The result is an identifier that is \emph{stable} (it does
not change when the content moves), \emph{universal} (it requires no
central authority), and \emph{self-authenticating} (anyone can verify it
by rehashing).

\simone{Reduce Git analogy repetition---appears four times across Sections 1, 3, 4, 6. Once is enough.}

Formal mathematics has not yet adopted this pattern.  Theorems in proof
assistants---Lean~\cite{lean4}, Coq~\cite{coq}, Isabelle~\cite{isabelle},
Agda~\cite{agda}---are identified by fully qualified names tied to
project structure, file paths, and namespace conventions.  Renaming a
file, moving a declaration to a different module, or reorganizing a
namespace hierarchy breaks every downstream reference, even though the
mathematical content is unchanged.  This problem is not hypothetical: it
occurs routinely in Mathlib, the largest Lean~4 library, which undergoes
continuous refactoring across its ${\sim}200{,}000$
declarations~\cite{mathlib4}.  A complementary line of work addresses
the downstream consequence---automated proof
repair~\cite{ringer2021repair, gopinathan2023repair}---but does not
eliminate the root cause: names are not intrinsic to the mathematics.

This paper asks: \emph{can we apply content-addressing to formal
mathematics?}  We show that the answer is yes, with important caveats
about the granularity of equivalence (Section~\ref{sec:equiv}).  We
describe a working implementation for Lean~4 and a decentralized
registry architecture that requires no infrastructure beyond Git.

\paragraph{Contributions.} A content-addressing pipeline for Lean~4 declarations:
canonicalization, serialization, and SHA-256 hashing (Section~\ref{sec:pipeline}).
Two canonicalization levels with different stability/equivalence trade-offs (Section~\ref{sec:canon}).
A Merkle DAG hashing mode that gives declarations truly content-determined identity across project boundaries (Section~\ref{sec:modes}).
A decentralized registry architecture exploiting the monotonicity of mathematical proof (Section~\ref{sec:registry}). Empirical evaluation on 57{,}613 theorems from a Mathlib
module environment (Section~\ref{sec:results}). An open-source implementation as a Lean~4 library and CLI tool (Section~\ref{sec:impl}). A working implementation
is available at \href{https://github.com/marcellop71/CA}{https://github.com/marcellop71/CA}

\moqian{Prose needs significant rewriting. Currently reads as AI-generated. Rewrite for clarity, academic voice, narrative coherence.}

% --------------------------------------------------------------------

% =====================================================================
\section{Related Work}
\label{sec:related}
% =====================================================================

\paragraph{Content-addressing in software and programming languages.}
Content-addressed storage is well-established in software systems
(Git~\cite{git}, IPFS~\cite{benet2014ipfs},
Nix~\cite{dolstra2004nix}).  The closest analog in programming
languages is Unison~\cite{unison}, which was designed from the ground
up around content-addressing: names are metadata, refactoring never
invalidates hashes, and the full syntax tree is hashed.  We differ in
three ways: we \emph{retrofit} content-addressing onto an existing
language not designed for it; we hash only the \emph{type}, not the
implementation (justified by proof irrelevance); and canonicalization
in a dependently typed setting is complicated by definitional equality
and universe polymorphism.
\simone{Deepen Unison comparison. Unison: designed around CA from start, hashes full AST, 512-bit SHA3. Ours: retrofit onto Lean, hash type only (proof irrelevance). Make structural difference explicit.}

\paragraph{Formal mathematics infrastructure.}
Hurd's OpenTheory~\cite{hurd2011opentheory} tackled the same core
problem---prover-independent theorem identity---for the HOL family,
using theory interpretations rather than cryptographic hashing.
\simone{Add OpenTheory (Joe Hurd, 2009--2014)---cross-prover package management for HOL family. Closest prior system. Doesn't use crypto hashing but same core problem. CRITICAL.}
Gauthier et al.~\cite{gauthier2019aligning} and M\"uller et
al.~\cite{muller2017alignment} align concepts \emph{across} type
theories; our content-addressing operates \emph{within} one, where
canonical forms can be compared syntactically---a natural prerequisite
for cross-system alignment.
\simone{Deepen alignment work discussion. Gauthier-Kaliszyk solve harder problem across foundations (HOL4/Coq/Mizar). Ours cleaner but narrower (single type theory). Explain intra-system CA as prerequisite for cross-system.}  The OMDoc/MMT
framework~\cite{rabe2013scalable, kohlhase2006omdoc} and the vision of
Kohlhase and Rabe~\cite{kohlhase2017qed} articulate a pluralistic
formal library; we provide a concrete identity mechanism for it.
Existing library management (Mathlib's monorepo~\cite{mathlib4}, Coq's
opam packages~\cite{coq}, Isabelle's AFP~\cite{isabelle_afp}) relies
on name-based identity throughout.

\sam{Discuss HashMath: independent CIC + content addressing + P2P. Complementary---HashMath from scratch, ours within Lean.}

\paragraph{Proof repair and deduplication.}
Ringer et al.~\cite{ringer2021repair} and Gopinathan et
al.~\cite{gopinathan2023repair} repair proofs broken by library
refactoring; content-addressing aims to prevent such breakage rather
than mitigate it.
\simone{Add proof repair literature: Ringer et al.\ PLDI 2021, Gopinathan et al.\ PLDI 2023. Complementary approach---automated repair vs stable identifiers.}  Type-based deduplication of theorem statements is
used in AI-for-theorem-proving pipelines~\cite{syntheticthm},
essentially Level~0 without cryptographic persistence.

\leo{Discuss Yatima project (github.com/argumentcomputer/yatima) as prior work. Dead but had nice ideas---content addressing + zkSNARK for type-checking.}

The question of
normal forms in dependent type theory~\cite{abel2018decidability,
coquand1991algorithm} underlies our canonicalization; we sidestep full
normalization with pragmatic partial levels.

% =====================================================================

% --------------------------------------------------------------------

\section{Name-based identity vs content-addressing}

This section describes the problem and the proposed solution at a
  high level.  We first show how Lean's name-based identity couples
  references to project structure, file layout, and naming conventions,
  making them fragile under refactoring and opaque across project
  boundaries.  We then introduce content-addressing---hashing the
  canonicalized type expression rather than the name---and state its
  key properties.

\paragraph{Name-based Identity}
% =====================================================================

To use a theorem in Lean~4, a programmer writes something like:

\begin{lstlisting}[caption={Using a theorem by its fully qualified name}]
import Mathlib.Data.Finset.Sum

example : forall (s : Finset Nat) (f g : Nat -> Int),
    s.sum (f + g) = s.sum f + s.sum g :=
  Finset.sum_add_distrib
\end{lstlisting}

This single reference, \texttt{Finset.sum\_add\_distrib}, encodes three
pieces of information beyond the mathematics: (1)~the \textbf{project}
(Mathlib), (2)~the \textbf{file}
(\texttt{Mathlib/Data/Finset/Sum.lean}, needed for the \texttt{import}),
and (3)~the \textbf{namespace path}
(\texttt{Finset.sum\_add\_distrib}, a human-chosen convention).
None of these are intrinsic to the mathematical statement.  The theorem
says that finite sums distribute over addition---a fact that is true
independently of what we call it or where we store it.

When any of these coordinates change, references break.  A theorem
renamed for consistency (\texttt{sum\_add\_distrib} $\to$
\texttt{sum\_add}), a declaration relocated to a different module
(\texttt{Data.Finset.Sum} $\to$ \texttt{Algebra.BigOps}), a namespace
hierarchy flattened or reorganized, a project forked with new
prefixes---each is a no-op mathematically, but a breaking change
syntactically.  The cost is borne by every downstream project: CI
pipelines fail, students' homework stops compiling, automated tooling
must track name migrations.

Within a single library like Mathlib, naming conventions and import
paths are at least centrally managed.  Across independent projects, the
situation is worse.  There is no standardized way for project~A to
refer to a result proved in project~B, other than adding a direct
dependency on B's entire package---including its specific version, file
layout, and naming scheme.  If two independent projects prove the same
theorem under different names, there is no mechanism to recognize that
they state the same mathematical fact.  Each project is an island, and
the only bridges are explicit package dependencies with name-level
coupling.

This problem will only grow.  Mathlib is evolving into the shared
foundation that virtually all Lean projects depend on, but the
ecosystem is expanding well beyond pure mathematics: formalization
projects in cryptography, blockchain protocol verification, physics,
and other applied domains are emerging, each building on Mathlib and
on each other.  As the number of independent projects grows, so does
the cost of name-level coupling---every cross-project reference is a
fragile link that any upstream refactoring can break.


% --------------------------------------------------------------------
\paragraph{Content-Addressing}
% --------------------------------------------------------------------

The central idea is simple: \emph{hash the mathematics, not the name.}

Given a declaration with a type expression (its mathematical statement),
we: (1)~\textbf{canonicalize} the type expression to remove
representation artifacts, (2)~\textbf{serialize} the canonical
expression to a byte sequence, and (3)~\textbf{hash} the byte sequence
with SHA-256.  The resulting 256-bit hash is the declaration's
\emph{content address}.
For theorems, the content address is therefore determined entirely by
the statement---the proof term is never hashed.
For propositions, this is inevitable given proof irrelevance.

Content addresses have four key properties.  They are \emph{stable}:
the address does not change when a declaration is renamed, moved, or
reorganized.  They are \emph{universal}: no project, file, or import
context is needed, and no coordination or central authority is required.
They are \emph{self-authenticating}: anyone can verify an address by
recomputing the hash from the declaration's type.  And they are
\emph{collision-resistant}: two mathematically different statements will
(with overwhelming probability) have different addresses.

An important caveat: content addresses identify declarations that share
the same \emph{canonical type}, but this does not make them freely
interchangeable.  Declarations carry metadata (attributes, instance
registrations) that the hash does not capture, and definitions with the
same type may compute differently.  We discuss these limits in
Section~\ref{sec:equiv}.

\subsection{Features}

  \paragraph{Search by statement, not by name}
  \label{rem:search}
  Content-addressing inverts the usual discovery workflow.  To determine
  whether a result is already known---proved, published as an open
  conjecture, or simply mentioned in a registry---one does not need to
  guess the name a library author might have chosen, nor search through
  documentation by keyword.  Instead, it suffices to \emph{write the
  statement} as a Lean type expression, compute its content address
  (a single call to the hashing pipeline), and look up that address
  in the registry.  A match means the exact mathematical fact is already
  recorded, regardless of the name, file, or project it lives in.
  A miss means no project in the registry has published that statement
  at the chosen canonicalization level.  The search is determined
  entirely by the mathematics---no text matching, no naming conventions,
  no knowledge of library structure required.

  \paragraph{Address is stable under proof improvement}
  [Stable under proof improvement]
  \label{rem:proof-improvement}
  Because the content address depends only on the statement, improving a
  proof---making it shorter, faster, or more elegant---does not change
  the address.  The mathematical ``API'' is stable even as
  implementations improve.  A library can silently replace a slow proof
  with a faster one and no downstream content address is affected.
  It's the analogue of "you can refactor the implementation without changing the interface" in software.

  \paragraph{Content addresses as permanent citations}
  \label{rem:citations}
  Academic papers about formal mathematics currently cite theorems by
  library name (e.g., \texttt{Finset.sum\_add\_distrib}).  These
  citations break silently when libraries refactor.  A content address
  serves as a permanent reference to a mathematical fact.
  A paper citing \texttt{ca:a1b2c3\ldots} remains
  valid regardless of future renaming or reorganization.

  \paragraph{Memoization for proof search (AI angle)}
  \label{rem:memoization}
  In automated or AI-assisted proof search, the same goal type often
  arises in different proof contexts.  A content address canonically
  identifies the goal independently of the surrounding proof state.
  This enables a global lookup table: before attempting synthesis, check
  whether the goal's content address already appears in the registry.
  If it does, the proof can be retrieved rather than re-discovered.

  \paragraph{Version-independent references}
  \label{rem:versions}
  A content address for a theorem is the same whether computed from
  Mathlib v4.3 or v4.9, provided the mathematical statement has not
  changed.  This decouples mathematical identity from library
  versioning: a project can depend on a \emph{mathematical fact} rather
  than on a specific version of a specific package.

  \paragraph{Progress tracking for formalization campaigns}
  \label{rem:campaigns}
  A formalization effort (e.g., all theorems in a textbook chapter) can
  be specified as a list of content addresses registered as open points.
  Progress is then the fraction of addresses whose status has moved from
  \texttt{open} to \texttt{proved}, a well-defined, cross-project
  metric that updates automatically via the status cascade.

  \paragraph{Toward granular building.}
  Lean's build system (Lake) operates at the granularity of entire
  modules: importing a single declaration requires compiling the full
  module and all its transitive dependencies.  This monolithic model
  compounds the cost of name-based coupling---a project that needs one
  lemma from a large module pays for thousands of unrelated declarations.
  Content-addressing opens a path toward finer granularity.  If
  declarations are individually addressable by content, a build system
  could in principle fetch and verify only the specific results a project
  actually uses, rather than pulling in entire module trees.  Granular
  addressing would lead naturally to granular imports: referencing
  individual facts rather than files.  This is a harder problem in a
  compiled ecosystem like Lean, where elaboration, compilation, and
  kernel-checking are tightly coupled at the module level, but
  content-addressing provides the necessary first step---a stable,
  module-independent name for each declaration.
  
% =====================================================================

% =====================================================================
\section{A working pipeline}
\label{sec:pipeline}
% =====================================================================

The address of a declaration is computed in three stages:
canonicalization, serialization, and hashing.

% --------------------------------------------------------------------
\subsection{Canonicalization}
\label{sec:canon}
% --------------------------------------------------------------------

Canonicalization transforms a Lean~4 \texttt{Expr} (expression) into a
canonical form where expressions differing only in irrelevant ways
become syntactically identical.  We implement two levels, offering a
trade-off between stability and the strength of the equivalence
relation.

\paragraph{Level~0: structural canonicalization.}
Level~0 is a pure transformation (no monadic context required) that
normalizes two sources of superficial variation.  First, \emph{universe
parameters}: Lean uses named universe polymorphism, and these names are
arbitrary; Level~0 renames them to positional indices (\texttt{u\_0},
\texttt{u\_1}, \ldots) by order of first occurrence.  Second,
\emph{metadata annotations}: Lean attaches \texttt{.mdata} nodes to
expressions for internal bookkeeping (source positions, elaboration
hints); these carry no mathematical content and are stripped.

Lean~4 uses de~Bruijn indices~\cite{debruijn1972} for bound variables,
so $\alpha$-equivalence is handled by the representation itself.
Binder \emph{names} (e.g., \texttt{n} vs.\ \texttt{x}) are cosmetic
and excluded from hashing.  Level~0 is fast and deterministic.  It
depends only on the structure of the \texttt{Expr} type, which is
stable within a major Lean version but has changed across major versions
(e.g., Lean~3 to Lean~4).  Its limitation is that it does not normalize
definitional equalities: two formulations that differ by an unfolded
abbreviation or a reduced beta-redex receive different hashes.

\paragraph{Level~1: reducible normalization.}
Level~1 extends Level~0 with \emph{reducible-transparency
normalization}.  It requires a \texttt{MetaM} (meta-programming monad)
context because it invokes Lean's \texttt{whnf} (weak head normal form)
reduction.  The normalizer, \texttt{deepNormalize}, recursively applies
\texttt{whnf} at every subexpression under \texttt{reducible}
transparency:

\begin{lstlisting}[style=pseudo,caption={Deep normalization (simplified pseudocode)}]
deepNormalize(e, depth):
  if depth >= 256: return e       -- recursion guard
  e <- whnf e                     -- reducible transparency
  match e:
    .app f a      -> .app (deepNormalize f) (deepNormalize a)
    .lam _ t b _  -> normalize t; open binder; normalize body;
                     re-abstract
    .forallE _ t b _ -> same as lambda
    .letE _ _ v b -> normalize v; substitute; normalize result
    .mdata _ e    -> deepNormalize e
    .proj _ _ e   -> normalize e; whnf the projection
    _             -> return e
\end{lstlisting}

Under reducible transparency, notation aliases (e.g.,
\texttt{HAdd.hAdd} $\to$ \texttt{Add.add}), declarations marked
\texttt{@[reducible]}, beta-redexes, let-bindings, and reducible
structure projections are all unfolded.  Crucially,
\texttt{@[semireducible]} declarations (including typeclass instances)
are \emph{not} unfolded, preserving abstraction boundaries:
\texttt{Add Nat} stays as \texttt{Add Nat} rather than being replaced
by its instance implementation.

As an example, consider \texttt{theorem add\_comm\_v1 : $\forall$ (a b
: Nat), a + b = b + a} and a variant \texttt{add\_comm\_v2 :
$\forall$ (a b : Nat), b + a = a + b}.  At Level~0, these may receive
different hashes if the elaborated expressions differ structurally.  At
Level~1, after \texttt{whnf} normalization, both reduce to the same
canonical form and receive the same hash.  On
\texttt{Init.Data.Nat.Basic}, the largest equivalence class grows from
391 members (Level~0) to 467 (Level~1)---76~additional declarations
collapsed by reducible unfolding.  Section~\ref{sec:results} presents a
larger-scale evaluation.

\franz{How useful is Level~1 in practice? Probably handles ``boring'' statements, but two mathematicians independently formalizing the same interesting theorem---would they get the same hash? Almost never at Level~1.}

\paragraph{Level~2: full definitional equality (open problem).}
The ideal canonicalization would assign the same hash to all
definitionally equal types.  However, Lean's dependent type theory does
not admit unique normal forms under full definitional
equality~\cite{abel2018decidability}.  Achieving Level~2 would require
either restricting the fragment of the type theory considered, or
developing new normalization techniques.  We leave this as an open
problem (see Section~\ref{sec:equiv} for further discussion).

\franz{Even Level~2 may not suffice---people use slightly different formalizations.}

\franz{Propose Level~3: formally different but ``spiritually the same''---one trivially proves the other. Quantify with fuzzy match + trivial proof detection (e.g.\ \texttt{exact?}\ or \texttt{apply?}\ in one step).}

% --------------------------------------------------------------------
\subsection{Serialization, Hashing, and Hashing Modes}
\label{sec:serial}
% --------------------------------------------------------------------

\paragraph{Serialization.}
The canonical expression is serialized to a \texttt{ByteArray} using a
tagged encoding: each \texttt{Expr} constructor is assigned a unique tag
byte, followed by the serialization of its children.  Binder names are
omitted (de~Bruijn indices suffice).  \texttt{BinderInfo} (implicit,
explicit, instance, etc.)\ is included for lambda and forall nodes.
\texttt{.mdata} is transparent: the child is serialized directly.  For
\texttt{.const name levels}, the treatment of \texttt{name} depends on
the hashing mode (below).

\sam{Hash-based proof compression: replace subexpressions by hash on network layer. Sam's ``hexpr'' constructor elides subterm structure for storage/communication.}

\paragraph{Hashing.}
The serialized byte array is hashed with SHA-256, producing a
64-character hexadecimal string.  We use SHA-256 via OpenSSL's EVP API
through a minimal C FFI shim, chosen for ubiquity and ease of binding
rather than performance (a faster hash such as BLAKE3 would work equally
well---the choice of hash function is orthogonal to the addressing
design).
\simone{Address: (a) mutual recursion in Merkle DAG, (b) why SHA-256 not BLAKE3, (c) literal \texttt{PROOF\_IRRELEVANT} sentinel value, (d) is L0 truly Lean-version-independent given \texttt{Expr} changes?}

\paragraph{Hashing modes.}
\label{sec:modes}
When serializing a constant reference (\texttt{.const name levels}),
there is a choice of what to emit for \texttt{name}.  In
\emph{name-based} mode (default), the declaration's \texttt{Name} is
serialized as a UTF-8 string---fast and stable within a single library,
but if a dependency is renamed, the address changes.  In
\emph{content-based} (Merkle DAG) mode, the 32-byte content hash of the
referenced declaration is emitted instead.  Declarations are processed
in topological (dependency) order via Kahn's
algorithm~\cite{kahn1962topological}, so each hash incorporates the
hashes of its transitive dependencies.  The result is a Merkle
DAG~\cite{merkle1988} where a declaration's identity depends only on the
mathematical content of everything it (transitively) references.

Mutually recursive declarations form cycles in the dependency graph.
The topological sort detects these and falls back to name-based hashing
for cycle members---a conservative choice that preserves soundness at
the cost of weaker identity for mutual definitions.  Mutual recursion
is uncommon in theorem \emph{statements}, but Lean's elaboration of
inductive types generates mutually recursive auxiliary declarations
(\texttt{rec}, \texttt{casesOn}, etc.), so the fallback affects more
declarations than one might initially expect when indexing definitions
alongside theorems.

\begin{table}[h]
\centering
\begin{tabular}{@{}lll@{}}
  \toprule
  \textbf{Mode} & \textbf{Constant references become} & \textbf{Trade-off} \\
  \midrule
  Name-based & The declaration's \texttt{Name} string & Fast; changes if deps renamed \\
  Content-based & 32-byte content hash of the dep & True content identity; slower \\
  \bottomrule
\end{tabular}
\caption{Hashing modes.}
\label{tab:modes}
\end{table}


% --------------------------------------------------------------------
\subsection{Results}
\label{sec:results}
% --------------------------------------------------------------------

This section presents both concrete examples of content-addressing in
action and a quantitative evaluation on Mathlib.

\paragraph{Duplicate detection.}
Large libraries inevitably accumulate duplicates---theorems with
different names but identical mathematical content.  Content-addressing
detects these automatically: duplicates share the same hash.  The
technique of comparing type expressions for deduplication is already
used in the AI-for-theorem-proving community---for instance, synthetic
theorem generation pipelines deduplicate candidates by \texttt{Expr}
equality~\cite{syntheticthm}.  Content-addressing provides a
principled, persistent, and cross-project version of what is currently
done ad~hoc within individual pipelines.

\paragraph{Cross-project references.}
Consider two independent projects that both work with a mathematical
concept:

\begin{lstlisting}[caption={Two projects proving the same lemma independently}]
-- In project-a/Algebra/Basic.lean
theorem add_comm_nat : forall a b : Nat, a + b = b + a :=
  Nat.add_comm

-- In project-b/MyMath/Nat.lean
theorem nat_addition_commutative : forall a b : Nat, a + b = b + a :=
  fun a b => Nat.add_comm a b
\end{lstlisting}

Despite different names, files, and projects, both declarations receive
the same content address.  A search tool can discover that these are the
same mathematical fact, enabling cross-project deduplication and
discovery without any coordination between the projects.

\paragraph{Open problems and conditional results.}
Content-addressing supports a structured workflow for open problems,
though this aspect of the system is prospective---the mechanism is
implemented but has not yet been exercised at scale.  A researcher
publishes a conjecture as an \emph{open point}---a \texttt{Prop} with a
content address but no proof:

\begin{lstlisting}[caption={An open conjecture and a conditional result}]
@[open_point "Bertrand's postulate, elementary form"]
def BertrandPostulate : Prop :=
  forall n : Nat, n >= 1 -> exists p, Nat.Prime p /\ n < p /\ p < 2 * n

@[publish "Consequence of Bertrand's postulate"]
theorem prime_between (n : Nat) (h : BertrandPostulate) (hn : n >= 1) :
    exists p, Nat.Prime p /\ n < p /\ p < 2 * n :=
  h n hn
\end{lstlisting}

The registry (Section~\ref{sec:registry}) tracks which conditional
results depend on which open points.  When someone eventually proves the
conjecture, the status of all downstream results cascades from
``conditional'' to ``proved.''  Several practical questions remain open:
how to handle conjectures that turn out to be false (or independent),
how to distinguish \texttt{sorry}-based development from genuinely
axiomatic assumptions, and how to curate status metadata across
projects.  We discuss these further in Section~\ref{sec:future}.
\simone{Mark open-problem workflows (Sections 5.3--5.4) as speculative/future work. No real usage evidence. Address: who curates status? What if open point is false? \texttt{sorry}-based vs axiomatic?}

\paragraph{Registry generation.}
A project publishes its annotated declarations by adding
\texttt{\#ca\_registry "registry/"} at the end of a Lean file that
imports the annotated modules.  When \texttt{lake build} compiles this
file, it writes \texttt{declarations.json} (address, name, status,
type, deps for each declaration) and \texttt{meta.json} (project
summary with counts by status) to the specified directory.

\paragraph{Evaluation on Mathlib.}
We evaluated the pipeline on the environment loaded by
\texttt{Mathlib.\allowbreak Algebra.\allowbreak Group.\allowbreak Basic}, a core module of Mathlib that
transitively imports 1{,}209 modules covering foundational arithmetic,
data structures, and algebraic hierarchy definitions.  The pipeline was
run in name-based mode at Level~1 (reducible normalization), restricted
to theorems.  All 57{,}613 theorems in the transitive environment were
processed successfully---no declarations were skipped or failed.

The 57{,}613 theorems produce 57{,}036 unique type addresses.  Of
these, 544 addresses are shared by more than one theorem, covering
1{,}121 theorems (1.9\% of the total).

\begin{table}[h]
\centering
\begin{tabular}{@{}rrr@{}}
  \toprule
  \textbf{Class size} & \textbf{Classes} & \textbf{Theorems} \\
  \midrule
  1 (unique)  & 56{,}492 & 56{,}492 \\
  2           & 513      & 1{,}026 \\
  3           & 29       & 87 \\
  4           & 2        & 8 \\
  \midrule
  Total       & 57{,}036 & 57{,}613 \\
  \bottomrule
\end{tabular}
\caption{Equivalence class size distribution (L1, name-based, theorems only).}
\label{tab:eval}
\end{table}

\paragraph{Examples of detected duplicates.}
The equivalence classes reveal genuine duplicates---theorems that exist
under multiple names due to historical renaming, namespace aliasing, or
independent re-statement:

\begin{lstlisting}[style=output,caption={Selected equivalence classes (L1)}]
Size 4:
  by_contra = by_contradiction
           = Classical.byContradiction = Classical.by_contradiction
  Type: (Not p -> False) -> p

Size 3:
  congr_arg = congrArg = let_val_congr
  Type: forall (b : a -> b), a = a' -> b a = b a'

  Nat.add_sub_cancel = Nat.add_sub_cancel_right
                     = Nat.add_sub_self_right
  Type: forall (a b : Nat), a + b - b = a

Size 2:
  mul_one = MulOneClass.mul_one
  Type: forall [MulOneClass M] (a : M), a * 1 = a

  mul_comm = CommMagma.mul_comm
  Type: forall [CommMagma G] (a b : G), a * b = b * a
\end{lstlisting}

The types shown are simplified pretty-printed representations of the
canonical forms; binder names and implicit arguments are artifacts of
the printer, not of the hash.
These duplicates arise from distinct sources.  The \texttt{by\_contra}
cluster contains four names for the same classical reasoning principle,
accumulated as Lean's standard library evolved.  The \texttt{mul\_one} /
\texttt{MulOneClass.mul\_one} pair reflects the algebraic hierarchy: the
root name and the typeclass-qualified name coexist as aliases.  The
\texttt{Nat.add\_sub\_cancel} triple shows the same arithmetic fact
named by three conventions.  In all cases, the content address correctly
identifies these as stating the same mathematical fact, modulo Level~1
canonicalization.

\paragraph{Coverage.}
All 57{,}613 theorems were hashed without error, demonstrating that
the pipeline (including the micro-batching workaround for
\texttt{whnf}; Appendix~\ref{app:engineering}) handles the full
transitive closure of a Mathlib module.  The current evaluation covers
a single module's environment at Level~1 only.  A systematic comparison
of Level~0 vs.\ Level~1 across multiple Mathlib entry points, and
evaluation in content-based (Merkle DAG) mode, remain as future work.

\simone{Quantitative evaluation across Mathlib: coverage, identical hashes at L0 vs L1, genuine duplicates, false positive rate, processing time. Produce a table across representative modules. CRITICAL.}



% =====================================================================

% =====================================================================
\section{Equivalence Classes and Their Limits}
\label{sec:equiv}
% =====================================================================

Two declarations sharing the same content address state the same
mathematical fact (modulo the canonicalization level).  The set of all
declarations with a given address forms an \emph{equivalence class}.

\paragraph{Same hash does not mean interchangeable.}  The content
address captures the \emph{type}, not the proof term, runtime behavior,
or elaboration metadata.  This has important consequences.
\emph{Theorems in direct proofs} are interchangeable, by proof
irrelevance.  \emph{Theorems in tactics} are not: a tactic like
\texttt{simp [Nat.add\_mul]} relies on the \texttt{@[simp]} attribute
attached to that specific \emph{name}; substituting an alias with the
same hash but different attributes would change the tactic's behavior.
\emph{Definitions} are not interchangeable: two definitions with the
same type may compute differently.  \emph{Instances} are not
interchangeable: typeclass instances are registered by name and
participate in instance resolution.

For this reason, resolution is always \emph{by name}.  The content
address is a stable, project-independent key for \emph{discovery} and
\emph{cross-referencing}, but resolution goes through a specific
declaration name that carries the necessary metadata.

\begin{lstlisting}[style=output,caption={Alias discovery via content address}]
Resolved: Nat.add_mul
  Hash:   e1ea8f392f15...  (L1)
  Module: Init.Data.Nat.Basic
  Kind:   theorem
  Type:   forall (n m k : Nat), (n + m) * k = n * k + m * k
  Aliases (1): Nat.right_distrib
\end{lstlisting}

\simone{Move limitations discussion earlier. Currently after practical examples that assume interchangeability. Foreshadow in Section~3.}

\paragraph{Canonicalization and definitional equality.}
The fundamental question underlying the equivalence relation is: which
expressions should receive the same hash?  The ideal answer is ``all
definitionally equal expressions,'' but this is not achievable in
general.  Lean's Calculus of Inductive Constructions (CIC) with proof
irrelevance does not admit unique normal forms under full definitional
equality~\cite{abel2018decidability, coquand1991algorithm}.  The
reduction rules---$\beta$, $\delta$ (at various transparencies),
$\iota$, $\zeta$, and proof irrelevance---do not confluently normalize
all expressions to a unique representative.

Our two-level approach is a pragmatic compromise.  Level~0 captures
syntactic equality modulo universe renaming and metadata---a
conservative but stable equivalence.  Level~1 captures equality modulo
reducible $\delta$-reduction, $\beta$-reduction, $\zeta$-reduction, and
reducible $\iota$-reduction---a strictly coarser equivalence that
collapses common aliases.  The gap between Level~1 and full definitional
equality is large: semireducible and opaque definitions are not
unfolded, so type synonyms at these transparencies create distinct
hashes.  In principle, one could add intermediate levels (semireducible
transparency, instance unfolding), but each level increases computation
cost, reduces stability across Lean versions, and risks divergent
normalization.

\paragraph{Universe polymorphism.}
Lean uses \emph{named} universe polymorphism: declarations are
parameterized over universe level variables with programmer-chosen
names.  Two declarations with identical bodies but different universe
parameter names (e.g., \texttt{.\{u, v\}} vs.\ \texttt{.\{a, b\}})
should hash identically.  Level~0 handles this by renaming universe
parameters to positional indices: the first universe variable
encountered becomes \texttt{u\_0}, the second \texttt{u\_1}, and so on.
The traversal order is deterministic (left-to-right depth-first over
the type expression).


% =====================================================================

% =====================================================================
\section{Implementation}
\label{sec:impl}
% =====================================================================

The implementation is a Lean~4 package (\texttt{CA}) available at
\url{https://github.com/marcellop71/CA} (repository to be made public
upon acceptance).  Downstream projects depend on it with a single line:

\begin{lstlisting}[style=pseudo]
require ca from git "https://github.com/marcellop71/CA" @ "main"
\end{lstlisting}

\paragraph{End-to-end workflow.}  A user who wants to content-address
their project performs four steps: (1)~add \texttt{require ca} to their
\texttt{lakefile.lean} and run \texttt{lake update}; (2)~annotate
declarations with \texttt{@[publish]} or \texttt{@[open\_point]}
(directly or retroactively via \texttt{attribute [publish] SomeName});
(3)~add \texttt{\#ca\_registry "registry/"} to a Lean file that
imports the annotated modules; (4)~run \texttt{lake build}---the
registry folder is generated automatically during compilation, no
separate tool invocation required.  The generated \texttt{registry/}
folder is committed to Git and pushed; the project is now a node in the
decentralized registry.

\paragraph{CLI.}  For batch processing (e.g., indexing all of Mathlib),
a CLI executable provides \texttt{ca address} (compute hashes, export
JSON manifest and TSV edge list) and \texttt{ca fetch} (store to
Redis).  Redis serves as an interactive lookup layer: given a content
address, it returns all known names, types, and modules in
constant time---a capability the static JSON manifest does not provide
efficiently at scale.

\begin{table}[h]
\centering
\begin{tabular}{@{}ll@{}}
  \toprule
  \textbf{Module} & \textbf{Role} \\
  \midrule
  \texttt{CA.Canonical}           & L0 (pure) and L1 (MetaM) canonicalization \\
  \texttt{CA.SHA256}              & SHA-256 via OpenSSL EVP (C FFI) \\
  \texttt{CA.ExprHash}            & Expr serialization, batch hashing, Merkle DAG \\
  \texttt{CA.Export}              & JSON manifest, TSV edge list, statistics \\
  \texttt{CA.Util}                & Dependency collection, declaration classification \\
  \texttt{CA.Registry.Basic}      & Entry types, persistent extensions, query API \\
  \texttt{CA.Registry.Attributes} & \texttt{@[publish]}, \texttt{@[open\_point]} registration \\
  \texttt{CA.Registry.Generate}   & \texttt{\#ca\_registry} command, JSON output \\
  \bottomrule
\end{tabular}
\caption{Module overview.}
\label{tab:modules}
\end{table}

\paragraph{Build and dependencies.}  The default \texttt{lake build}
compiles the \texttt{CA} library, which requires only OpenSSL.  The CLI
executable (\texttt{lake build ca}) additionally requires hiredis for
Redis-backed lookups.  A Nix flake provides a development shell with the
Lean toolchain and all native dependencies for both Linux and macOS.
The implementation currently targets Lean~4.29.0-rc1 and tracks the
Lean release cycle; it depends on the Batteries
library~\cite{batteries} for \texttt{NameMapExtension}.

\simone{Add end-to-end user walkthrough. Current implementation section is half a page.}


% =====================================================================

% =====================================================================
\subsection{Limitations and Future Work}
\label{sec:limits}
% =====================================================================

\paragraph{Current limitations.}
Level~0 and Level~1 canonicalization do not capture full definitional
equality; the gap between Level~1 and the ideal is large and may be
fundamentally unclosable.  Level~1 addresses depend on Lean's
\texttt{whnf} behavior, which can change across toolchain versions;
Level~0 addresses are stable within a major Lean version, while Level~1
should be considered best-effort within a given toolchain.  Lean's
native \texttt{whnf} (implemented in C) lacks heartbeat checks and
crashes on large environments; a micro-batching workaround is effective
but inelegant (Appendix~\ref{app:engineering}).  Applying \texttt{whnf}
to definition values causes process-level memory growth that crashes
regardless of batch boundaries, limiting content-addressing of
definitions (as opposed to theorems, where proof irrelevance makes value
hashing unnecessary).  Finally, Lean~4's elaboration pipeline does not
support user-defined fallbacks for unresolved names; our resolution
mechanism requires explicit syntax (\texttt{resolve\%}), preventing
seamless integration.

\paragraph{Future work.}
\phantomsection\label{sec:future}
Several directions are worth pursuing.  \emph{Resolution elaborators}:
implement \texttt{use}, \texttt{use!}, and \texttt{resolve\%} as Lean~4
elaborators that query the registry, resolve by content address, and
dynamically import modules.  The intended syntax (not yet implemented):

\begin{lstlisting}[caption={Prospective resolution syntax (not yet implemented)}]
-- resolve% looks up a content address in the registry and
-- returns the name of a matching declaration in scope
example : forall a b : Nat, a + b = b + a :=
  resolve% "a1b2c3d4e5f6..."   -- resolves to Nat.add_comm
\end{lstlisting}

\emph{Ecosystem-wide indexing}: apply
content-addressing to the entirety of Mathlib to produce a global index
of equivalence classes, enabling large-scale duplicate detection,
library health metrics, and structured navigation.
\emph{Open-problem workflows}: develop the open-point and
conditional-result mechanism (Section~\ref{sec:results}) into a robust
workflow, addressing false conjectures, \texttt{sorry}-based vs.\
axiomatic assumptions, and cross-project status curation.
\emph{AI-assisted proof search}: content addresses provide a stable,
type-based index for machine learning models to retrieve relevant lemmas
by querying the content address of a goal type.
\emph{Cross-system alignment}: explore content-addressing in other proof
assistants (Coq, Isabelle, Agda) and investigate whether addresses can
be made comparable via shared canonicalization.
\emph{Intermediate canonicalization levels}: investigate canonicalization
at semireducible transparency or with selective instance unfolding.
\emph{Formal verification}: prove (in Lean itself) that the
canonicalization and serialization stages preserve the intended
equivalence relation.

\moqian{CA + ZKP reduces verification from O(compile entire dependency chain) to O(1). Critical for AI-generated proofs at scale.}

\moqian{CA shortens proof writing: hash as stable ID, readable name as alias. DNS model---IP is identity, domain name is human-friendly.}

\paragraph{Engineering Obstacles}
\label{app:engineering}

A practical obstacle to Level~1 canonicalization at scale: Lean's native
\texttt{whnf} (implemented in C) contains no heartbeat or memory
checks.  When processing large environments (${\sim}58{,}000$
declarations in Mathlib modules), two failure modes arise.  First,
\emph{MetaM state accumulation}: a single \texttt{MetaM} session
crashes after ${\sim}25{,}000$ \texttt{whnf} calls due to cumulative
growth of caches, discriminant trees, and instance tables; the crash
point depends on call count, not on which specific declarations are
processed.  Second, \emph{process-level memory growth}: applying
\texttt{whnf} to definition values causes monotonic RSS growth that
eventually crashes regardless of batch boundaries.

\paragraph{Micro-batching.}  Our workaround is to process declarations
in micro-batches of~10, each in a fresh \texttt{MetaM} session.  Types
receive Level~1 normalization; values are always canonicalized with
Level~0 only (no \texttt{whnf}).  For theorems, the value hash is the
sentinel \texttt{PROOF\_IRRELEVANT} (proof irrelevance).  The result is
an asymmetric strategy: \textbf{types get L1, values get L0}.  Hash
computation is split into two phases (canonicalize+hash vs.\
pretty-print) in separate \texttt{MetaM} sessions, with a flush every
5{,}000 declarations to bound live \texttt{Expr} objects in memory.

An upstream fix---heartbeat checks and resource guards in Lean's native
\texttt{whnf} implementation---would eliminate the need for
micro-batching entirely.

% =====================================================================

% =====================================================================
\section{Decentralized Registry}
\label{sec:registry}
% =====================================================================

Content addresses give every declaration a stable,
location-independent identity.  The registry organizes the mutable
metadata---who proved what, what is open, what depends on what---around
this immutable core.

Every formalization in Lean already lives in a Git repository.  The
decentralized registry exploits this directly: each participating
project hosts a \texttt{registry/} folder in its own repo, and the
global registry is the union of all such folders.  No additional
infrastructure is needed---no IPFS, no blockchain, no central server.

\sam{Compare P2P DHT (Sam's approach) vs Git-based registry (ours). Sam found DHT doesn't scale, moved to local-first + DHT for publishing locations only.}

This simplicity is possible because \textbf{mathematical consensus is
trivial}.  Unlike distributed systems that must agree on ordering or
resolve write conflicts, a mathematical registry has a monotonic truth
condition: a proposition is proved if \emph{any} repository contains a
valid proof.  Proofs cannot be retracted---once proved, always proved.
Two repositories proving the same proposition under different names is
not a conflict; it is redundancy, and the content address shows they
state the same fact.  There is nothing to vote on, no quorum to reach.
This observation connects to the broader vision of a pluralistic formal
library where multiple communities contribute to a shared body of
verified knowledge~\cite{kohlhase2017qed}.

In concrete terms:

\begin{itemize}[nosep]
  \item \textbf{Publication} is \texttt{git push}---a project generates
    its \texttt{registry/} folder, commits it, and pushes.
  \item \textbf{Verification} is \texttt{lake build}---Lean's kernel
    type-checks the proofs; anyone can clone and verify independently.
  \item \textbf{Consensus} is trivial: one valid proof is enough.
  \item \textbf{Discovery} is a curated list of known repositories
    (\texttt{sources.json}).
\end{itemize}

\paragraph{Annotation attributes.}
Authors annotate their code with two attributes provided by the
\texttt{CA.Registry} module.
\texttt{@[open\_point]} marks a \texttt{def X : Prop} as an open
problem---a mathematical statement published without a proof.
Applying it to a theorem or a non-\texttt{Prop} definition is a
compile-time error.
\texttt{@[publish]} marks any declaration (theorem, definition, axiom)
for publication to the registry.
Both attributes accept an optional description string and support
retroactive application from a separate file, allowing annotation of
existing code or upstream dependencies without modifying the original
source:

\begin{lstlisting}[caption={Retroactive annotation of existing declarations}]
import CA
import MyProject.Theorems

attribute [open_point "My conjecture"] MyProject.SomeConjecture
attribute [publish "Key lemma"]        MyProject.some_lemma
\end{lstlisting}

\paragraph{Declaration status.}
Status is computed automatically from the Lean environment:

\begin{table}[h]
\centering
\begin{tabular}{@{}lll@{}}
  \toprule
  \textbf{Pattern} & \textbf{Status} & \textbf{Rule} \\
  \midrule
  \texttt{theorem foo : P := proof} & proved & Type-checks, no \texttt{sorry} \\
  \texttt{def P : Prop} with \texttt{@[open\_point]} & open & Explicitly marked \\
  \texttt{theorem bar (h : OpenP) : Q} & conditional & Proved modulo open hypotheses \\
  \bottomrule
\end{tabular}
\caption{Automatic status classification.}
\label{tab:status}
\end{table}

For conditional declarations, the registry records which open addresses
are required---the declarations whose proof would promote the entry
from conditional to proved.

\paragraph{Registry folder.}
The \texttt{registry/} folder is generated locally inside each project.
The \texttt{\#ca\_registry "registry/"} command (or the \texttt{ca
registry} CLI) loads the compiled environment, collects annotated
declarations, computes content addresses, classifies status, and writes
two files:

\begin{itemize}[nosep]
  \item \texttt{declarations.json}---the append-only manifest.  Each
    entry contains the content address, name, module, kind, status,
    pretty-printed type, and type dependencies.
  \item \texttt{meta.json}---project-level metadata (hash level,
    open/proved/conditional counts).
\end{itemize}

Entries are never removed.  Status updates (open $\to$ proved) are
reflected by updating the status field in place; Git history provides
the audit trail.

\paragraph{Status cascade.}
When a proof is submitted for an open point, conditional declarations
downstream may become unconditionally proved:

\begin{enumerate}[nosep]
  \item A previously open address acquires a proof---status becomes
    \texttt{proved}.
  \item All \texttt{conditional} entries whose open dependencies
    included the newly proved address are re-evaluated.
  \item If all of an entry's open dependencies are now proved, it is
    promoted to \texttt{proved}.
  \item This cascades transitively.
\end{enumerate}

Cross-project cascades emerge when a consumer merges multiple
registries: an open point in project~A may be proved in project~B,
promoting conditional entries in project~C, with no coordination
between A, B, and~C.  The dependency graph also makes \emph{impact}
computable: the impact of an open point is the number of conditional
declarations transitively downstream of it.

\paragraph{Sources and discovery.}
Because the registry is fully decentralized, a discovery mechanism is
needed so that resolution tools know where to look.  A curated
\texttt{sources.json} file in the CA repository enumerates known Git
repositories that host a \texttt{registry/} folder:

\begin{lstlisting}[style=json,caption={\texttt{sources.json} format}]
{ "registries": [
    { "url": "https://github.com/leanprover-community/mathlib4",
      "registry_path": "registry/" },
    { "url": "https://github.com/user/my-project",
      "registry_path": "registry/" }
  ]
}
\end{lstlisting}

Adding a project to the ecosystem requires five steps: (1)~add
\texttt{require ca} to the project's \texttt{lakefile.lean};
(2)~annotate declarations with \texttt{@[publish]} or
\texttt{@[open\_point]}; (3)~add \texttt{\#ca\_registry "registry/"}
to a Lean file that imports the annotated modules;
(4)~\texttt{lake build}, commit \texttt{registry/}, and push;
(5)~open a PR adding one entry to \texttt{sources.json}.
Step~5 is the only coordination point; everything else is local to the
project.  The static list is the right starting point---it can later be
replaced by automatic discovery (e.g., a GitHub topic tag or a DHT)
without changing the rest of the architecture.

\paragraph{Current maturity.}  The architecture described here is
implemented but early-stage.  Discovery relies on a manually curated
\texttt{sources.json}, there is no mechanism for revoking or correcting
erroneous entries, and metadata conflicts (e.g., two projects assigning
different status to the same address) are resolved by the consumer at
merge time.  The theoretical elegance of monotonic consensus is real,
but the tooling around it is minimal.


% =====================================================================

% =====================================================================
\section{Conclusion}
\label{sec:conclusion}
% =====================================================================

We have presented a content-addressing system for formal mathematics,
implemented as a Lean~4 library.  By hashing canonicalized type
expressions, we assign declarations stable identifiers determined by
their mathematical content rather than their human-chosen names.

Content-addressing does not require full normalization to be useful.
Even partial canonicalization (Level~0) provides meaningful stability
and cross-project discoverability.  The decentralized registry
architecture, made possible by the monotonicity of mathematical proof,
requires no infrastructure beyond Git---though significant engineering
work remains to make it practical at ecosystem scale.

The system is a working prototype, and important challenges
remain---most fundamentally, the gap between our canonicalization levels
and full definitional equality.  We hope this work demonstrates that
content-addressing is a productive lens through which to view the
organization of formal mathematical knowledge, and that the tools
presented here are a useful step toward a more interconnected ecosystem
for formalized mathematics.

\simone{Separate language-agnostic parts (registry format, discovery protocol, Merkle DAG) from Lean-specific parts (canonicalization, \texttt{Expr} serialization, \texttt{whnf} batching). Former could become standard for other proof assistants.}


% =====================================================================
% References
% =====================================================================



\bibliographystyle{plain}
\bibliography{references}

\end{document}
